%%%%%%%%%%%%%%%%%%%%%%%%%%%%%%%%%%%%%%%%%%%%%%%%%%%%%%%%%%%%%%%%%%%%%%%%%%%%%%%
% Differential entropy and continuous random variables
%%%%%%%%%%%%%%%%%%%%%%%%%%%%%%%%%%%%%%%%%%%%%%%%%%%%%%%%%%%%%%%%%%%%%%%%%%%%%%%

\subsection{Why a New Notion Is Needed}

For a discrete random variable, entropy is maximized by the uniform distribution and is always nonnegative.
If we try to naively extend
\(H(X)=-\sum_x p(x)\log p(x)\)
by replacing the sum with an integral, we obtain \emph{differential entropy}, which behaves quite differently.

A useful way to remember what changes is to think about \emph{quantization}.
Let $X$ be continuous and define a quantized version $X^\Delta$ by binning with resolution $\Delta>0$.
Then $H(X^\Delta)$ is a meaningful number of bits to describe the quantized value.
As $\Delta\to 0$, $H(X^\Delta)$ diverges, so there cannot be a finite ``absolute'' entropy for an exact real number.
Differential entropy captures the \emph{finite part} of this divergence.

\subsection{Differential Entropy}
\label{sec:diffent-def}

\begin{definition}[Differential entropy]
Let $X$ be a continuous random variable with pdf $f_X$ on $\fR^d$.
Its differential entropy is
\begin{equation}
    h(X) \coloneqq -\int_{\fR^d} f_X(x)\,\log f_X(x)\,dx.
\end{equation}
\end{definition}

\paragraph{Discretization identity (heuristic).}
For scalar $X$ quantized at resolution $\Delta$,
\begin{equation}
    H(X^\Delta) = h(X) + \log\frac{1}{\Delta} + o(1), \qquad \Delta\to 0.
\end{equation}
So $H(X^\Delta)\to\infty$ as $\Delta\to 0$.

\paragraph{Key caveats.}
Differential entropy:
\begin{itemize}[leftmargin=2em]
    \item can be \textbf{negative};
    \item is \textbf{not invariant} under one-to-one transformations;
    \item does \textbf{not} directly correspond to a lossless code length (because exact reals cannot be losslessly encoded with finite bits).
\end{itemize}

\begin{example}[Negative differential entropy]
If $X\sim\mathrm{Unif}(0,\tfrac12)$, then $f_X(x)=2$ on $(0,\tfrac12)$ and
\[
    h(X) = -\int_0^{1/2} 2\log 2\,dx = -\log 2 < 0.
\]
\end{example}

\begin{theorem}[Scaling rule]
If $X$ is scalar and $a\neq 0$, then
\begin{equation}
    h(aX)=h(X)+\log|a|.
\end{equation}
More generally, if $X\in\fR^d$ and $A\in\fR^{d\times d}$ is invertible,
\(h(AX)=h(X)+\log|\det A|\).
\end{theorem}


\subsubsection*{Joint and Conditional Differential Entropy}

For a pair $(X_1,X_2)$ with joint density $f_{X_1,X_2}$, define
\[
h(X_1,X_2)
=-\iint f_{X_1,X_2}(x_1,x_2)\,\log f_{X_1,X_2}(x_1,x_2)\,\mathrm{d}x_1\mathrm{d}x_2,
\]
and the conditional entropy
\[
h(X_1\mid X_2)
=-\iint f_{X_1,X_2}(x_1,x_2)\,\log f_{X_1\mid X_2}(x_1\mid x_2)\,\mathrm{d}x_1\mathrm{d}x_2.
\]
These satisfy the chain rule:
\[
h(X_1,X_2)=h(X_2)+h(X_1\mid X_2).
\]

\subsection{KL Divergence and Mutual Information for Continuous Variables}
\label{sec:diffent-mi}

The quantities that remain well-behaved in the continuous setting are KL divergence and mutual information.

\begin{definition}[KL divergence for densities]
Let $f$ and $g$ be pdfs on $\fR^d$.
Define
\begin{equation}
    \KL{f}{g} \coloneqq \int f(x)\,\log\frac{f(x)}{g(x)}\,dx,
\end{equation}
with the convention that if $f(x)>0$ but $g(x)=0$ on a set of positive $f$-measure, then $\KL{f}{g}=+\infty$.
\end{definition}

\begin{definition}[Mutual information for continuous random variables]
Let $(X,Y)$ have joint density $f_{X,Y}$ with marginals $f_X$ and $f_Y$.
The mutual information is
\begin{equation}
    I(X;Y)
    \coloneqq \KL{f_{X,Y}}{f_X f_Y}
    = \E{\log\frac{f_{X,Y}(X,Y)}{f_X(X)f_Y(Y)}}.
\end{equation}
\end{definition}

All properties from the discrete case that are based on KL divergence still hold:
\begin{itemize}[leftmargin=2em]
    \item $I(X;Y)\ge 0$ with equality iff $X$ and $Y$ are independent.
    \item $I(X;Y)=I(Y;X)$.
    \item Data processing inequality (Theorem~\ref{thm:dpi}).
\end{itemize}

\paragraph{Chain rule with differential entropy.}
Whenever the relevant entropies exist,
\begin{equation}
    I(X;Y)=h(X)-h(X\mid Y)=h(Y)-h(Y\mid X).
\end{equation}
Although $h(\cdot)$ itself shifts under invertible transformations, these shifts cancel in $I(X;Y)$.

\begin{remark}[Invariance of mutual information]
If $a\neq 0$ and $X$ is scalar, then
\[
I(aX;Y)=h(aX)-h(aX\mid Y)=(h(X)+\log|a|)-(h(X\mid Y)+\log|a|)=I(X;Y).
\]
More generally, mutual information is invariant under smooth bijections applied separately to $X$ and $Y$.
\end{remark}

\subsection{Gaussian Maximizes Differential Entropy Under a Variance Constraint}
\label{sec:diffent-gauss-max}

Just as the uniform distribution maximizes discrete entropy on a finite set, the Gaussian maximizes differential entropy under a second-moment constraint.

\begin{theorem}[Maximum differential entropy]
\label{thm:gauss-maxent}
Let $X$ be scalar with $\E{X}=0$ and $\E{X^2}\le P$.
Then
\begin{equation}
    h(X) \le \frac12\log\bigl(2\pi e P\bigr),
\end{equation}
with equality if and only if $X\sim\cN(0,P)$.
\end{theorem}

\begin{proof}
Let $g$ be the pdf of $\cN(0,P)$ and let $f$ be the pdf of $X$.
KL non-negativity gives
\[
0\le \KL{f}{g} = \int f(x)\log\frac{f(x)}{g(x)}\,dx = -h(X) - \int f(x)\log g(x)\,dx.
\]
For the Gaussian,
\(\log g(x) = -\tfrac12\log(2\pi P) - \tfrac{x^2}{2P}\log(e)\)
(because our $\log$ is base $2$).
Thus
\[
\int f(x)\log g(x)\,dx
= -\tfrac12\log(2\pi P) - \tfrac{\E{X^2}}{2P}\log(e)
\ge -\tfrac12\log(2\pi P) - \tfrac12\log(e)
= -\tfrac12\log(2\pi e P),
\]
where we used $\E{X^2}\le P$.
Substituting back yields $h(X)\le \tfrac12\log(2\pi e P)$.
Equality forces $\E{X^2}=P$ and $\KL{f}{g}=0$, i.e., $f=g$.
\end{proof}
